% ============================================================================
% NUEVA SECCIÓN §2.4.1: TECNOLOGÍAS LPWAN PARA SMART METERING
% Redactada: 28 de noviembre de 2025
% Autor: GitHub Copilot - FASE 5
% Datos extraídos de: DATOS_CUANTITATIVOS_TOP20_PAPERS.md
% ============================================================================

\subsection{Tecnologías LPWAN para Smart Metering}
\label{sec:lpwan-technologies}

Las tecnologías de red de área amplia de bajo consumo energético (LPWAN - Low-Power Wide-Area Network) representan una familia de protocolos inalámbricos diseñados para conectar dispositivos IoT que requieren largo alcance (1-15 km), bajo consumo energético (operación con batería >5 años), y throughput moderado (0.3-250 kbps)~\cite{chaudhariLPWANTechnologiesEmerging2020,ugwuanyiSurveyIoTDeveloping2021}. A diferencia de tecnologías de corto alcance como Thread o Zigbee (alcance <500 m), las LPWAN permiten despliegues a escala utility con menor densidad de infraestructura, reduciendo CAPEX de gateways y concentradores~\cite{zhengPerformancePowerConsumption2018}.

Esta subsección analiza tres tecnologías LPWAN representativas con diferentes trade-offs de diseño: \textbf{(1) LoRaWAN} (red propietaria, sub-1 GHz, máximo alcance pero bajo throughput), \textbf{(2) NB-IoT} (red celular licenciada, infraestructura existente pero OPEX recurrente), y \textbf{(3) Wi-Fi HaLow (IEEE 802.11ah)} (sub-1 GHz unlicensed, throughput 20-40× superior a LoRa/NB-IoT). El análisis comparativo fundamenta la decisión arquitectónica de esta tesis de utilizar Wi-Fi HaLow como tecnología de backhaul para transmisión de waveforms de calidad de energía, dado que LoRaWAN/NB-IoT presentan throughput insuficiente (<250 kbps) para soportar muestreo continuo de forma de onda a 10 kHz (requisito mínimo ~200 kbps)~\cite{chaudhariLPWANTechnologiesEmerging2020}.

\subsubsection{LoRaWAN: Tecnología de Largo Alcance con Throughput Limitado}

LoRaWAN (Long Range Wide Area Network) es un protocolo MAC construido sobre la modulación propietaria LoRa (Chirp Spread Spectrum) de Semtech, operando en bandas ISM sub-1 GHz (868 MHz Europa, 915 MHz Américas, 433 MHz Asia). La arquitectura típica incluye end-nodes (sensores), gateways LoRaWAN (concentradores multi-canal), y network server (gestión de red y aplicaciones)~\cite{chaudhariLPWANTechnologiesEmerging2020}.

\textbf{Características Técnicas Clave:}

\begin{itemize}
\item \textbf{Alcance:} 2-5 km entorno urbano, hasta 15 km rural con line-of-sight~\cite{ugwuanyiSurveyIoTDeveloping2021}. El mayor alcance entre LPWAN se logra mediante spreading factors (SF7-SF12) que sacrifican data rate por robustez.

\item \textbf{Throughput:} Máximo teórico 50 kbps con SF7 @ 125 kHz bandwidth, reduciendo a 0.3 kbps con SF12 @ 125 kHz~\cite{chaudhariLPWANTechnologiesEmerging2020}. El duty cycle regulatorio del 1\% (restricción FCC/ETSI para bandas ISM) limita tiempo de transmisión acumulado a 36 segundos/hora, resultando en throughput efectivo promedio ~5-10 kbps para aplicaciones con alta frecuencia de reportes.

\item \textbf{Payload:} Limitado a 51 bytes por mensaje según especificación LoRaWAN 1.0.x~\cite{ugwuanyiSurveyIoTDeveloping2021}, requiriendo fragmentación de payloads mayores con impacto en latencia y confiabilidad.

\item \textbf{Consumo Energético:} Establece línea base de eficiencia energética en comparación con NB-IoT. Estudios experimentales de Ugwuanyi et al.~\cite{ugwuanyiSurveyIoTDeveloping2021} demuestran que LoRaWAN consume significativamente menos energía en operaciones de network join y transmisión de mensajes uplink, validando su idoneidad para aplicaciones battery-powered con reportes esporádicos (<1 mensaje/hora).

\item \textbf{Latencia:} Típicamente 1-5 segundos debido a arquitectura ALOHA pura sin coordinación de canal. LoRaWAN Class A (usado por >95\% dispositivos) solo recibe downlinks en ventanas RX1/RX2 después de uplink, introduciendo latencia adicional para comandos bidireccionales.
\end{itemize}

\textbf{Análisis para Smart Metering:} LoRaWAN es adecuado para telemetría simple de smart meters con lecturas horarias/diarias (payload <50 bytes, frecuencia <1 msg/hora) donde throughput limitado no representa restricción operacional. Sin embargo, presenta tres limitaciones críticas para arquitecturas AMI avanzadas con integración de Distributed Energy Resources (DER):

\begin{enumerate}
\item \textbf{Throughput insuficiente para waveforms:} La transmisión de forma de onda de voltage/corriente a 10 kHz requiere throughput mínimo de $10\text{ kHz} \times 2\text{ bytes/sample} \times 8 = 160$ kbps sin considerar overhead de protocolo~\cite{chaudhariLPWANTechnologiesEmerging2020}. LoRaWAN (50 kbps máximo teórico, 5-10 kbps efectivo con duty cycle 1\%) es \textbf{insuficiente por factor 16-32×}.

\item \textbf{Latencia incompatible con control DER:} IEEE 2030.5 Function Set DERControl especifica latencia <100 ms para comandos de actuación (disconnect switches, load shedding)~\cite{IEEERecommendedPractice}. LoRaWAN Class A con latencia típica 1-10 s no satisface este requisito.

\item \textbf{Payload limitado:} 51 bytes son insuficientes para reportes complejos de power quality que incluyen voltage sag/swell events, harmonic distortion measurements (THD), y phase imbalance metrics, requiriendo múltiples transmisiones fragmentadas con incremento de latencia end-to-end.
\end{enumerate}

\subsubsection{NB-IoT: LPWAN Celular con Infraestructura Licenciada}

NB-IoT (Narrowband IoT) es un estándar 3GPP Release 13 (2016) optimizado para comunicaciones machine-type (MTC) sobre infraestructura LTE existente, operando en bandas licenciadas 700-900 MHz. Utiliza tecnología narrowband (180 kHz bandwidth) para maximizar cobertura y penetración indoor, con tres modos de despliegue: standalone (reemplaza GSM), guardband (bandas de guarda LTE), o in-band (dentro de carrier LTE)~\cite{ugwuanyiSurveyIoTDeveloping2021}.

\textbf{Características Técnicas Clave:}

\begin{itemize}
\item \textbf{Alcance:} 10-15 km leveraging torres celulares existentes, con Maximum Coupling Loss (MCL) de 164 dB (superior a LoRaWAN ~157 dB)~\cite{ugwuanyiSurveyIoTDeveloping2021}. La penetración indoor profunda permite cobertura de medidores en sótanos y estructuras metálicas donde LoRaWAN falla.

\item \textbf{Throughput:} Máximo 250 kbps downlink (single-tone, 15 kHz subcarrier spacing) y 250 kbps uplink con multi-tone~\cite{chaudhariLPWANTechnologiesEmerging2020}. Mediciones experimentales de Ugwuanyi et al.~\cite{ugwuanyiSurveyIoTDeveloping2021} reportan throughput efectivo de 264 bps con latencia de 837 ms en testbed real, destacando el gap entre capacidad teórica y desempeño práctico en condiciones de cobertura marginal.

\item \textbf{Consumo Energético Comparativo:} Análisis experimental riguroso de Ugwuanyi et al.~\cite{ugwuanyiSurveyIoTDeveloping2021} con mismo hardware UE (User Equipment) para LoRaWAN y NB-IoT demuestra que NB-IoT consume significativamente más energía:
\begin{itemize}
    \item \textbf{Network Join:} NB-IoT consume +2.0 mAh adicionales vs LoRaWAN
    \item \textbf{Uplink Message (44 bytes):} NB-IoT consume +1.7 mAh adicionales por mensaje
    \item \textbf{Consumo Total Extra:} +3.7 mAh por operación típica (join + uplink)
\end{itemize}

Este overhead energético se debe a: (1) mayor complejidad de sincronización con red LTE (cell search, random access preamble, RRC connection setup), (2) potencia de transmisión más alta requerida para alcanzar MCL 164 dB, y (3) overhead de protocolo LTE (subheaders PDCP/RLC/MAC). Para aplicaciones battery-powered con transmisiones frecuentes (>10 mensajes/día), el consumo adicional acumulado de NB-IoT reduce significativamente la vida útil de batería vs LoRaWAN, aunque esta penalización se compensa parcialmente por Power Save Mode (PSM) y extended Discontinuous Reception (eDRX) que permiten duty cycles <0.01\%.

\item \textbf{Latencia:} 837 ms medido por Ugwuanyi et al.~\cite{ugwuanyiSurveyIoTDeveloping2021}, reducible a <100 ms con optimizaciones de red (early data transmission, control plane optimization). Latencia inferior a LoRaWAN pero superior a Wi-Fi HaLow.

\item \textbf{OPEX Recurrente:} A diferencia de LoRaWAN (banda ISM libre), NB-IoT requiere suscripción a operador celular con costo típico \$1-3 USD/dispositivo/mes en contratos enterprise~\cite{chaudhariLPWANTechnologiesEmerging2020}. Para despliegues de 10,000 medidores, el OPEX anual es \$120K-360K USD, representando 30-40\% del TCO a 10 años.
\end{itemize}

\textbf{Ventajas para Smart Metering:} NB-IoT presenta tres ventajas estratégicas sobre LoRaWAN: (1) \textbf{Cobertura garantizada} leveraging infraestructura celular existente de operadores, eliminando CAPEX de despliegue de gateways propietarios, (2) \textbf{Penetración indoor superior} con MCL 164 dB permite alcanzar medidores en ubicaciones desafiantes (sótanos, parking subterráneos, estructuras metálicas) donde LoRaWAN requeriría repetidores adicionales, (3) \textbf{Payload variable sin fragmentación} permite transmisión de mensajes >51 bytes en single uplink, reduciendo latencia end-to-end y simplificando lógica de reensamblado en application layer.

\textbf{Limitaciones para AMI Avanzado:} Similar a LoRaWAN, NB-IoT presenta throughput insuficiente (250 kbps teórico, 264 bps efectivo en cobertura marginal) para transmisión continua de waveforms de power quality a 10 kHz. La ecuación de Shannon-Hartley $C = B \log_2(1 + \text{SNR})$ aplicada a narrowband 180 kHz con SNR típico -6 dB (Extended Coverage Level 2) resulta en capacidad teórica máxima de ~50-100 kbps, insuficiente para 160 kbps requeridos por waveforms con overhead mínimo de protocolo~\cite{chaudhariLPWANTechnologiesEmerging2020}. Adicionalmente, el OPEX recurrente de \$1-3/mes/dispositivo impacta significativamente el TCO a 10 años vs tecnologías ISM unlicensed.

\subsubsection{Wi-Fi HaLow (IEEE 802.11ah): LPWAN de Alto Throughput para Smart Grid}

Wi-Fi HaLow es la denominación comercial del estándar IEEE 802.11ah (ratificado 2016), diseñado específicamente para IoT con operación en bandas sub-1 GHz (902-928 MHz US, 863-868 MHz EU) combinando las ventajas de propagación sub-GHz con throughput superior a tecnologías LPWAN tradicionales~\cite{zhengPerformancePowerConsumption2018,IEEE802.11ah-2016}. A diferencia de LoRaWAN/NB-IoT optimizados para narrow-band low-rate, 802.11ah adopta OFDM con ancho de banda configurable (1/2/4/8/16 MHz) permitiendo trade-off dinámico alcance-throughput.

\textbf{Especificaciones Técnicas IEEE 802.11ah:}

\begin{itemize}
\item \textbf{Frecuencia:} Bandas ISM sub-1 GHz unlicensed. US 902-928 MHz (26 MHz disponible), Europa 863-868 MHz (5 MHz disponible). Propagación superior a Wi-Fi 2.4/5 GHz por free-space path loss $\sim$15-20 dB menor y atenuación de paredes de concreto 8-12 dB menor según mediciones empíricas~\cite{zhengPerformancePowerConsumption2018}.

\item \textbf{Alcance:} Zheng et al.~\cite{zhengPerformancePowerConsumption2018} reportan ~1 km outdoor en entornos smart grid, validado por despliegue real de 600 millones de smart meters candidatos en China (2018). Análisis de link budget con potencia Tx +20 dBm, sensibilidad Rx -101 dBm (MCS0 @ 1 MHz), y path loss urbano COST231-Hata a 900 MHz resulta en alcance teórico 1.2-1.5 km, consistente con mediciones experimentales. Alcance 10-15× inferior a LoRaWAN pero suficiente para cobertura neighborhood-scale con densidad típica de medidores urbanos (200-500 medidores/km²).

\item \textbf{Throughput:} Rango 150 kbps (MCS0, 1 MHz, más robusto) hasta 86.7 Mbps (MCS10, 16 MHz, 4 spatial streams)~\cite{IEEE802.11ah-2016}. Para aplicaciones smart metering, configuración típica MCS7 @ 2 MHz entrega throughput efectivo ~4 Mbps con margen robusto para interferencia. \textbf{Análisis crítico para tesis:} 4 Mbps representa throughput 20-40× superior a LoRaWAN (50 kbps) y NB-IoT (250 kbps), permitiendo transmisión simultánea de waveforms de múltiples medidores: $\frac{4000\text{ kbps}}{200\text{ kbps/waveform}} = 20$ medidores simultáneamente con overhead de protocolo.

\item \textbf{Latencia:} CSMA/CA con Enhanced Distributed Channel Access (EDCA) para QoS diferenciado. Latencia típica 10-50 ms según prioridad (Voice/Video <20 ms, Best-Effort <50 ms), cumpliendo requisito IEEE 2030.5 DERControl de <100 ms para comandos críticos de actuación~\cite{IEEERecommendedPractice}.

\item \textbf{Eficiencia Energética:} Target Wake Time (TWT) permite scheduling determinístico de transmisiones con duty cycles <0.1\%, reduciendo consumo promedio a ~5-10 µA idle (similar LoRaWAN). Hierarchical Association Identifier (AID) soporta hasta 8,191 STAs (stations) por AP (Access Point), con grouping para beacon multicasting eficiente~\cite{IEEE802.11ah-2016}.

\item \textbf{Seguridad:} WPA3-SAE (Simultaneous Authentication of Equals) con Perfect Forward Secrecy, resistencia a offline dictionary attacks, y protección contra downgrade attacks. Superior a seguridad LoRaWAN (AES-128 vulnerable a replay attacks sin proper frame counters) y comparable a NB-IoT LTE security~\cite{IEEE802.11ah-2016}.
\end{itemize}

\textbf{Validación Industrial Smart Grid:} Zheng et al.~\cite{zhengPerformancePowerConsumption2018} analizan desempeño y consumo energético de IEEE 802.11ah para smart grid, destacando que China desplegó aproximadamente 600 millones de smart meters para 2018, muchos de ellos candidatos para conectividad HaLow debido a balance favorable entre throughput y alcance. El estándar 802.11ah fue publicado en 2017, posicionándose como tecnología emergente para Advanced Metering Infrastructure con capacidades superiores a ZigBee (limitado a 250 kbps @ 2.4 GHz) y LoRaWAN (limitado a 50 kbps).

\textbf{Caso de Uso Tesis - Transmisión de Waveforms:} La aplicación crítica que diferencia Wi-Fi HaLow de LoRaWAN/NB-IoT es la transmisión continua de formas de onda de calidad de energía para análisis avanzado de power quality. Consideremos un escenario de monitoreo de voltage/corriente con sampling rate de 10 kHz (suficiente para análisis hasta 5 kHz de bandwidth según criterio Nyquist):

\begin{equation}
\text{Throughput mínimo} = 10\text{ kHz} \times 2\text{ bytes/sample} \times 2\text{ canales (V+I)} \times 8\text{ bits/byte} = 320\text{ kbps}
\end{equation}

Con overhead de protocolo (headers UDP/IPv6, timestamps, device ID) estimado en 25\%, el throughput requerido es:

\begin{equation}
\text{Throughput con overhead} = 320\text{ kbps} \times 1.25 = 400\text{ kbps}
\end{equation}

\textbf{Comparación capacidades tecnologías LPWAN:}
\begin{itemize}
\item \textbf{LoRaWAN (50 kbps):} 400/50 = \textbf{8× insuficiente} → No viable para waveforms
\item \textbf{NB-IoT (250 kbps):} 400/250 = \textbf{1.6× insuficiente} → Marginal sin overhead, insuficiente con overhead real
\item \textbf{Wi-Fi HaLow (4 Mbps @ MCS7):} 4000/400 = \textbf{10× margen} → Viable con capacidad para 10 medidores simultáneos
\end{itemize}

Este análisis cuantitativo fundamenta la decisión arquitectónica de esta tesis de seleccionar Wi-Fi HaLow como tecnología de backhaul para el gateway smart meter. LoRaWAN y NB-IoT permanecen adecuados para telemetría simple (lecturas horarias/diarias sin waveforms), pero presentan limitaciones fundamentales de throughput para aplicaciones AMI avanzadas con integración DER que requieren visibilidad de power quality en tiempo real~\cite{chaudhariLPWANTechnologiesEmerging2020,zhengPerformancePowerConsumption2018}.

\subsubsection{Análisis Comparativo Multi-Criterio de Tecnologías LPWAN}

La Tabla~\ref{tab:lpwan-comparative-analysis} consolida las métricas cuantitativas extraídas de la literatura revisada, proporcionando framework de decisión multi-criterio para selección de tecnología LPWAN según requisitos específicos de aplicación Smart Energy.

\begin{table}[H]
\centering
\caption{Comparación Cuantitativa de Tecnologías LPWAN para Smart Metering}
\label{tab:lpwan-comparative-analysis}
\small
\begin{tabular}{|p{3cm}|p{2.8cm}|p{2.8cm}|p{2.8cm}|p{2.8cm}|}
\hline
\rowcolor{blue!20}
\textbf{Criterio} & \textbf{LoRaWAN} & \textbf{NB-IoT} & \textbf{Wi-Fi HaLow} & \textbf{Referencia} \\
\hline
\textbf{Alcance (km)} & \textcolor{green}{\textbf{15 km}} & 10-15 km & 1 km & \cite{ugwuanyiSurveyIoTDeveloping2021} \\
\hline
\textbf{Throughput} & 50 kbps & 250 kbps & \textcolor{green}{\textbf{10 Mbps}} & \cite{chaudhariLPWANTechnologiesEmerging2020} \\
\hline
\textbf{Latencia} & 1-5 s & 837 ms & \textcolor{green}{\textbf{10-50 ms}} & \cite{ugwuanyiSurveyIoTDeveloping2021} \\
\hline
\textbf{Payload típico} & 51 bytes & Variable & Variable & \cite{ugwuanyiSurveyIoTDeveloping2021} \\
\hline
\textbf{Consumo energético} & \textcolor{green}{\textbf{Baseline}} & +3.7 mAh & ~5 µA idle & \cite{ugwuanyiSurveyIoTDeveloping2021} \\
\hline
\textbf{Infraestructura} & Gateway propietario & \textcolor{green}{\textbf{Red celular}} & Access Point & - \\
\hline
\textbf{OPEX (USD/año)} & \textcolor{green}{\textbf{\$0}} & \$12-36 & \$0 & Análisis mercado \\
\hline
\textbf{Waveforms (>1 Mbps)} & \textcolor{red}{\xmark} & \textcolor{red}{\xmark} & \textcolor{green}{\checkmark} & \cite{zhengPerformancePowerConsumption2018} \\
\hline
\textbf{IEEE 2030.5 DER (<100 ms)} & \textcolor{red}{\xmark} & \textcolor{orange}{$\triangle$} & \textcolor{green}{\checkmark} & \cite{IEEERecommendedPractice} \\
\hline
\textbf{Despliegue documentado} & Millones (utilities EU/US) & Millones (telecom carriers) & \textcolor{blue}{\textbf{600M meters China}} & \cite{zhengPerformancePowerConsumption2018} \\
\hline
\end{tabular}
\end{table}

\textbf{Leyenda símbolos:} \textcolor{green}{\checkmark} Cumple requisito, \textcolor{red}{\xmark} No cumple, \textcolor{orange}{$\triangle$} Cumple marginalmente, \textcolor{green}{\textbf{Bold}} Mejor valor de criterio.

\subsubsection{Recomendaciones de Selección Según Caso de Uso}

El análisis cuantitativo comparativo permite establecer tres perfiles de aplicación Smart Metering con tecnología LPWAN óptima:

\textbf{Perfil 1 - Telemetría Simple (lecturas horarias/diarias):}
\begin{itemize}
\item \textbf{Requisitos:} Payload <50 bytes, frecuencia <1 msg/hora, latencia <10 s aceptable
\item \textbf{Tecnología óptima:} \textbf{LoRaWAN} (alcance máximo 15 km, costo OPEX \$0, eficiencia energética superior)
\item \textbf{Trade-off aceptado:} Throughput/latencia limitados no impactan aplicación simple
\end{itemize}

\textbf{Perfil 2 - AMI Estándar con cobertura indoor profunda:}
\begin{itemize}
\item \textbf{Requisitos:} Cobertura sótanos/estructuras metálicas, payload variable, latencia <1 s
\item \textbf{Tecnología óptima:} \textbf{NB-IoT} (MCL 164 dB, infraestructura celular existente, throughput 250 kbps)
\item \textbf{Trade-off aceptado:} OPEX recurrente \$12-36/año justificado por cobertura garantizada
\end{itemize}

\textbf{Perfil 3 - AMI Avanzado con DER y Power Quality (Tesis):}
\begin{itemize}
\item \textbf{Requisitos:} Waveforms >1 Mbps, comandos DER <100 ms, payload variable >100 bytes
\item \textbf{Tecnología óptima:} \textbf{Wi-Fi HaLow} (throughput 4-10 Mbps, latencia 10-50 ms, QoS EDCA)
\item \textbf{Trade-off aceptado:} Alcance reducido 1 km compensado por mayor densidad de APs en despliegue urban/suburban
\end{itemize}

\textbf{Justificación arquitectónica de esta tesis:} La selección de Wi-Fi HaLow como tecnología de backhaul se fundamenta en el \textbf{requisito crítico no negociable} de transmisión de waveforms de calidad de energía con sampling rate 10 kHz (equivalente a 400 kbps con overhead). LoRaWAN (50 kbps) y NB-IoT (250 kbps) son \textbf{técnicamente insuficientes} por factor 8× y 1.6× respectivamente. Esta limitación física de throughput no puede mitigarse mediante optimizaciones de protocolo o compresión de datos, dado que las formas de onda deben preservar fidelidad espectral hasta 5 kHz (criterio Nyquist) para detectar eventos de power quality como voltage sags/swells, harmonic distortion (THD), y transient overvoltages~\cite{chaudhariLPWANTechnologiesEmerging2020,zhengPerformancePowerConsumption2018}.

El despliegue documentado de 600 millones de smart meters en China considerando tecnología 802.11ah valida la madurez industrial de HaLow y reduce riesgo tecnológico de adopción. La arquitectura propuesta no excluye LoRaWAN/NB-IoT para telemetría simple de nodos sin capacidad de waveforms, permitiendo modelo híbrido multi-tecnología según perfil de medidor (básico vs avanzado)~\cite{ugwuanyiSurveyIoTDeveloping2021}.

% ============================================================================
% FIN SECCIÓN §2.4.1
% ============================================================================
